\section{Conclusion: Peripherality as a Trajectory}
\label{sec:conclusion}

The flat-world promise was half right. Falling costs of moving goods, data
and people have brought the edge closer in absolute terms: more islanders are
online, connected and richer than ever. But integration has also turned
remoteness into a recurring tax, paid in freight, bandwidth, talent and
negotiating capacity, and the taxes reinforce one another.

The data in this paper support that picture in some places and complicate it
in others. Distance, more than smallness, carries an income penalty. Shipping
connectivity at the smallest, most remote nodes has stagnated while the
network thickened elsewhere. Redundancy scales with size, and a single cable
still carries entire national economies. Small islands lose most of their
graduates. Technologies now arrive quickly but are used far less intensively.
Against this, internet access has converged dramatically, gaps with the old
shipping hubs have narrowed, and four peripheral economies show that the edge
can be escaped---though the richest of them turn out to be small rather than
far.

The right conclusion is neither fatalism nor complacency. Peripherality is a
trajectory, not a location. Whether a place moves towards the centre or away
from it depends on how it pools scale, buys redundancy, finances its
exposure, bargains with others and keeps its people connected. In an
accelerating world, the cost of standing still at the edge rises every year;
so does the return to acting early.

\begin{pullquote}
The edge is not a place. It is a direction of travel.
\end{pullquote}
